% Le manuscrit BnF NAL 1644 et l'édition par Anderson (1615) du
% « De aequationum recognitione et emendatione » de Viète.
% Compiler avec : tectonic -X compile article.tex
\documentclass[11pt,a4paper]{article}
\usepackage{iftex}
\ifPDFTeX
  \usepackage[T1]{fontenc}
  \usepackage[utf8]{inputenc}
  \usepackage{lmodern}
\else
  \usepackage{fontspec}
\fi
\usepackage[french]{babel}
\usepackage{amsmath,amssymb,amsthm}
\usepackage{array,booktabs}
\setlength{\emergencystretch}{2em}
\usepackage[margin=2.8cm]{geometry}
\usepackage[hidelinks]{hyperref}

\newtheorem{proposition}{Proposition}
\theoremstyle{remark}
\newtheorem*{remarque}{Remarque}

% Lien vers une vue de Gallica du manuscrit.
\newcommand{\vue}[1]{\href{https://gallica.bnf.fr/ark:/12148/btv1b10033620q/f#1.item}{vue~#1}}
\newcommand{\vues}[2]{\href{https://gallica.bnf.fr/ark:/12148/btv1b10033620q/f#1.item}{vues~#1}--#2}

\title{Le manuscrit BnF NAL 1644 et l'exemplaire d'Anderson :\\
une collation du \emph{De aequationum recognitione et emendatione}\\
de Viète (Paris, 1615)}
\author{Michel Hua\\\small Chercheur indépendant}
\date{Novembre 2026}

\begin{document}
\maketitle

\begin{abstract}
En 1615, Alexander Anderson publie les deux traités de Viète sur les
équations d'après un exemplaire qu'il dit «~corrompu en bien des endroits,
mutilé en quelques-uns~», et dresse à la fin de sa préface la liste de ce
qu'il y a suppléé ou corrigé. Quatre articles de cette liste désignent
expressément l'exemplaire : trois passages qui y manquaient, une phrase qu'il
a récrite. Le manuscrit latin NAL 1644 de la Bibliothèque nationale de
France, qui contient les deux mêmes traités, présente aux quatre endroits le
texte que décrit Anderson, et porte dans sa marge, au point d'insertion où
l'imprimé place l'addition, le passage qu'il a imprimé, dans des termes
voisins des siens. Quatre fautes que sa liste corrige sont aussi dans le texte
du manuscrit, dont la marge corrige deux. Le lien entre le manuscrit et
l'édition d'Anderson n'est pas nouveau : à une question de H.~Braid (1902),
Paul Tannery répondait qu'il avait collationné le manuscrit avec les
\emph{Opera} de 1646, qu'il y voyait une copie ancienne qui représentait, si
l'on croit le récit d'Anderson, sa rédaction définitive, et que quelques
corrections marginales lui paraissaient de la main d'Anderson ; il ne
désignait aucun passage où se verrait le travail de l'éditeur. Faite sur
l'édition de 1615 et sur la liste d'Anderson, la collation précise cette
lecture : aux quatre endroits, le texte du manuscrit a l'état de
l'exemplaire, et la marge porte la rédaction d'Anderson. Quelques faits
s'y accordent mal, et la collation ne permet pas de choisir entre trois
hypothèses : le manuscrit serait l'exemplaire d'Anderson, une copie de cet
exemplaire, ou un témoin de la même source corrigé d'après l'imprimé. Les lectures du manuscrit
viennent d'une transcription automatique qu'aucun humain n'a encore
vérifiée sur les images.
\end{abstract}

\section{L'édition de 1615 et l'exemplaire d'Anderson}\label{sec:anderson}

Viète meurt en 1603. Le \emph{De aequationum recognitione et emendatione
tractatus duo} paraît douze ans plus tard à Paris, chez Jean Laquehay, par
les soins du mathématicien écossais Alexander Anderson
\cite{anderson1615}. Le volume réunit le \emph{De recognitione aequationum}
(p.~1--64) et le \emph{De emendatione aequationum} (à partir de la p.~65) ;
les \emph{Opera mathematica} de 1646, éditées par Frans van Schooten,
reprennent les deux textes aux p.~84--158 \cite{opera1646}.

La préface «~Ad Matheseos studiosos~» dit d'où vient le texte. Anderson
remercie «~Iacobo Alelmo~», Jacques Aleaume, ingénieur du roi, «~qui quidem
Vietae aduersaria liberrime mecum communicauit, ex quibus haec sunt
deprompta~» : les papiers de travail de Viète, dont ces traités sont tirés.
Puis il décrit son propre travail sur l'exemplaire dont il s'est servi,
«~non paucis in locis depravato, quibusdam etiam mutilo~»\,: chaque équation
et chaque preuve numérique des exemples ont été réexaminées, les fautes
corrigées, et ce qui paraissait manquer au texte a été «~restitué selon
l'exigence du sens~» et soumis au jugement du lecteur
\cite[préface, 2\textsuperscript{e} page]{anderson1615}.

Ce jugement, il le rend possible par une liste de deux pages placée après la
préface et intitulée «~Quae nos ad plenam Autoris mentem in exemplari deesse
visa suppleuimus, aut a Typographis \& Anagnostis commissa siue omissa
animaduertimus, haec fere sunt~» \cite[liste, après la préface]{anderson1615}.
Le titre mêle deux catégories : ce qui manquait dans l'exemplaire et a été
suppléé, et ce que les typographes et les correcteurs ont commis ou omis.
Pour la plupart des trente-deux articles, la liste ne dit pas à quelle
catégorie ils appartiennent. Quatre seulement nomment l'exemplaire :

\begin{enumerate}
\item p.~58, \emph{De recognitione}, chap.~18, prop.~14 : «~deerant in
Exemplari haec~», et suit la définition de $Z$ \emph{plano-solidum} ;
\item p.~60, chap.~19, prop.~4 : «~deerant in exemplari nostro~», et suit
«~differentia est B, \& fit Z plano-planum ex ductu utriusuis extremae, in
cubum differentiae reliquarum, alterne sumptarum~» ;
\item p.~98, \emph{De emendatione}, «~Theor.~2. ad probl.~2.~» : «~haec
verba exemplaris, \emph{ex applicatione quadrati aggregati laterum, minus
triplo rectangulo sub lateribus, ad rectangulum sub lateribus}, ita
mutauimus \emph{ex applicatione cubi a D plano, ad differentiam quadruplam
quadratorum a D \& B}~» ;
\item p.~99, «~Theor.~3. proxime sequenti~» : «~haec prorsus deerant in
exemplari~», et suit «~\& fit Z plano-planum ex applicatione cubi a D plano,
ad aggregatum quadruplum quadratorum a D \& B~».
\end{enumerate}

\noindent Dans l'imprimé, les trois passages suppléés et la phrase récrite
sont dans le corps du texte, aux p.~58, 60, 98 et~99 ; la liste en garde la
trace. Elle se clôt sur «~De his aliisque penes lectorem iudicium esto~».

On sait par ailleurs, grâce aux lettres que cite Bigourdan, que les papiers
de Viète étaient passés à Aleaume, et qu'à sa mort, en 1627, ses livres et
ses papiers furent dispersés : Peiresc le déplore le 11~novembre 1627, et
Jacques Dupuy lui répond que «~quelques-uns des héritiers~» se sont
accommodés des ouvrages de Viète \cite[p.~237--238]{bigourdan1916}. Le
prénom n'est pas sûr : de Thou, cité par Henry, dit les papiers de Viète
confiés par ses héritiers à «~Petri Alealmi Aurelianensis~», Pierre
Aleaume d'Orléans, là où Anderson nomme Jacques \cite{henry1880}. Ce que
devint l'exemplaire d'Anderson n'est pas connu des sources consultées ici.

\section{Le manuscrit NAL 1644}\label{sec:manuscrit}

Le volume latin des Nouvelles acquisitions 1644 vient du fonds Libri
(Ashburnham, Libri~1201). Delisle le décrit en 1888 : «~Papiers de François
Viète. Volume jadis coté G~», contenant les deux traités, puis «~Fragments
astronomiques~» \cite[p.~166, n\textsuperscript{o}~CXX]{delisle1888} ;
l'inventaire des nouvelles acquisitions de 1891 enregistre ensemble les
papiers de Viète des NAL 1643 et~1644 \cite[t.~II, p.~664]{delisle1891}. Il est
numérisé dans Gallica \cite{gallica} ; on le cite ici par les vues de cette
numérisation, et par la série de numéros à l'encre portée en haut des recto
(1 à~79), qui court sans lacune à travers tout le volume et que l'écrivain
cite lui-même.

Une couverture (vue~1) précède trois petits feuillets lettrés «~A~»,
«~B~» et «~C~» (vues~2--4), d'une écriture serrée que la transcription dit
autre que celle des traités. Le feuillet~A (\vue{2}) traite, par la
«~duplicata hypostasis~» et la «~climactica paraplerosis~», des équations
qui s'écriraient aujourd'hui $x^6 + px^3 = q$ et $x^6 + px^2 = q$, puis,
sous le titre «~Facillima resolutio aequationum cubicarum affectarum~»,
les ramène par un changement d'inconnue à une équation cubique «~a radice
simplici~» ; le
feuillet~B (\vue{3}) énonce un théorème sur les différences des termes pris
alternativement dans une suite de proportionnelles ; le feuillet~C est
blanc. C'est sans doute le «~court fragment sur la solution de l'équation
du troisième degré~» que Tannery place au commencement du manuscrit et
qu'il attribue à Anderson \cite{tannery1902} ; la transcription n'en a pas
comparé l'écriture à celle des marges. Le
\emph{De recognitione} occupe les \vues{5}{68}, sous le titre «~Francisci
Vietae De Recognitione Aequationum tractatus~», en vingt et un chapitres ; le
\emph{De emendatione} les \vues{69}{141}, sous le titre «~Francisci Vietae. De
aequationum Emendatione tractatus secundus~», en quatorze chapitres, jusqu'à
«~Finis~». Les chapitres sont ceux de l'imprimé, dans le même ordre ; chaque
exemple numérique des feuillets est dans l'imprimé. Les \vues{142}{163}
(feuillets~67 à~79 de la série à l'encre) sont de l'astronomie sans nom
d'auteur ; ce sont les feuillets que Pantin cite comme «~copie partielle~»
de l'\emph{Harmonicon coeleste} de Viète \cite[p.~123,
n.~22]{pantin2003}.

Les feuillets sont de tailles inégales : de petits feuillets sont reliés entre
de plus grands et en recouvrent une partie sur les vues. Le texte des traités
est d'une même cursive posée. Les marges sont pleines de corrections : un mot
ou un nombre souligné dans la ligne, le bon écrit dans la marge entre deux
traits. Une main plus rapide, à l'encre noire, annote et corrige à partir de
la vue~13 ; elle écrit «~a nobis~». Aux \vues{61}{66}, des notes et calculs
d'une écriture plus petite occupent la marge, sans que la transcription
puisse dire si c'est une autre main ; aux \vues{105}{119}, les marges
paraissent de la main et de l'encre du texte. Bigourdan, qui cherchait le
manuscrit de l'\emph{Harmonicon}, a comparé les écritures : «~la majeure
partie du n\textsuperscript{o}~1644 des Nouv. acq. lat., jusqu'au
feuillet~66, pourrait être de la même main que le n\textsuperscript{o}~7274
du fonds latin~», volume qu'il tient pour le brouillon de l'\emph{Harmonicon}
et peut-être pour un autographe de Viète, tout en notant que la comparaison
ne fournit «~pas de résultat absolument probant~»
\cite[p.~238--240]{bigourdan1916}.

Les avis sur les mains divergent. Tannery tenait le manuscrit pour «~une
copie ancienne (pas plus de la main d'Anderson que de celle de Viète)~»
\cite{tannery1902}, déjà décrite en 1896 comme «~copie ancienne~»
\cite[p.~209--210]{tannery1896} ; Bigourdan, on l'a vu, envisage pour les
feuillets~1 à~66 la main de Viète. Pour les feuillets d'astronomie, le
catalogue de Van~Egmond range les ff.~67r--79v parmi des «~miscellaneous
autograph pages~» \cite{vanegmond1985}, et Busard les donnerait aussi
pour autographes \cite{busard1964}, ce qui n'est connu ici que de
seconde main ; Tannery y voyait «~un brouillon de Viète~» pour un
théorème de l'\emph{Harmonicon}, Pantin une «~copie partielle~».

Une main plus tardive renvoie à l'édition de 1646 sur la couverture et aux
\vue{5}, \vue{69} et \vue{141} («~p.~127. editionis 1646.~», «~pag.~158
editionis elzevir. anni 1646~»). Bigourdan signale sur la couverture et les
marges «~quelques notes de la main de Lalande, qui l'avait comparé~» aux
\emph{Opera} de 1646 \cite[p.~240]{bigourdan1916} ; ce sont vraisemblablement
celles-ci.

\subsection*{La question de Braid et la réponse de Tannery}

En 1902, H.~Braid demande dans l'\emph{Intermédiaire des mathématiciens}
ce qui, dans le traité de 1615, revient à Viète et ce qui revient à
Anderson, et si «~le manuscrit de la Bibliothèque Nationale de Paris
[nouv. acq. lat., 1644 (Libri, 1201)]~» peut éclairer la question
\cite{braid1902}. Tannery répond dans le même volume. Il a collationné
«~le manuscrit Bibl. Nat. lat. n. a. 1644, avec l'édition de 1646~» : les
différences sont «~exclusivement d'ordre typographique~», et portent
surtout sur les notations, moins abrégées dans le manuscrit. La copie est
«~passablement incorrecte~», et plusieurs fautes grossières se retrouvent
dans l'imprimé ; il cite «~\emph{vocum} au lieu de \emph{vitium}~»
(1646, p.~134, l.~12 en remontant). «~Si l'on regarde comme vrai le récit
d'Anderson~», conclut-il, cette copie représente sa rédaction définitive,
d'après les brouillons de Viète. Il ajoute que les marges portent, «~de
plusieurs mains différentes~», des corrections qui «~n'ont pas passé dans
le texte imprimé~», et que «~quelques-unes paraissent de l'écriture
d'Anderson~», connue par ses notes sur l'exemplaire du Diophante de Bachet
conservé à la Sorbonne \cite[réimpr. p.~410--411]{tannery1902}. Il avait
déjà relevé en 1896 les différences de notation et confronté une équation
de la «~page 13~v.~» du manuscrit à la p.~98 de l'édition Elzevier
\cite[p.~209--210]{tannery1896}.

Le rapprochement entre NAL 1644 et le travail d'Anderson est donc fait
depuis 1902, mais sans le détail. Tannery a collationné le manuscrit sur
les \emph{Opera} de 1646, non sur l'édition de 1615 ni sur la liste où
Anderson dit ce qu'il a suppléé, et il n'indique aucun passage où se
verrait sa rédaction. Ces passages existent, et ils déplacent un peu sa
conclusion : aux quatre endroits que la liste attribue à l'exemplaire, le
texte du manuscrit a l'état d'avant Anderson, et sa rédaction est dans la
marge.

Trois notes marginales, toutes de lecture douteuse, peuvent indiquer que le
texte est une copie : «~in exemplari excripto / vel potius adscriptum~»
(vue~20, lue en partie), «~transcripta ex exemplari~» en tête d'une addition
encadrée (\vue{100}, lecture probable) et «~deerant haec in exemplari~»,
biffé (\vue{114}, lecture très incertaine). Cette dernière ne correspond à
aucun article de la liste d'Anderson.

\section{La collation}\label{sec:collation}

\subsection{Les quatre passages}

\paragraph{1. \emph{De recognitione}, chap.~18, prop.~14 (\vue{62}).}
L'énoncé est
$B\,A^2 - A^5 = Z$ («~Si B solidum in A quadratum $-$ A quadratocubo
aequetur Z planosolido~»). Le texte du manuscrit définit $B$ à l'aide de
trois proportionnelles, puis passe à «~vel a cubo compositae ex secunda et
tertia~\dots~» sans avoir défini $Z$. Un «~1\textsuperscript{o}~» est écrit
au-dessus de ce «~vel~», et la marge, de la main plus petite, porte au même
appel une note qui définit $Z$ : «~et [\dots] Z planosolidum quod fit
[\dots] solido [\dots] cubo a duabus [\dots] in quadratum compositae ex
prima et secunda~». Elle est biffée en plusieurs endroits et tachée ; la
transcription la donne pour «~très incomplète~». Ce qui se lit est compatible
avec la phrase d'Anderson, «~\& Z plano-solidum quod fit a B solido minus
cubo a duabus primis, in quadratum compositae ex prima \& secunda~», et
l'imprimé de la p.~58 insère cette phrase exactement devant «~vel a cubo
compositae ex secunda \& tertia~», là où le manuscrit place son appel.

\paragraph{2. \emph{De recognitione}, chap.~19, prop.~4 (\vue{64}).}
Le texte du manuscrit s'arrête après «~Sunt quatuor continue proportionales
quarum alterne sumptarum~» sur une croix, sans verbe, et reprend «~dum
intelligitur prima, minor inter extremas~\dots~». À la même croix, la marge
donne : «~differentia est B et fit Z planoplanum ex ductu utriusque [lecture
douteuse] inter extremas in cubum differentiae reliquarum alterne
sumptarum~», suivi d'une vérification sur l'exemple $1, 2, 4, 8$. L'imprimé
de la p.~60 porte, au même endroit : «~quarum alterne sumptarum differentia
est B : \& fit z plano-planum, ductu utriusuis extremae, in cubum
differentiae reliquarum alterne sumptarum, \& dum intelligitur prima~\dots~».
La proposition suivante (prop.~5, \vue{65}), de même forme, a dans le texte
du manuscrit la clause entière, «~differentia est B. fit autem Z
solidosolidum ductu utriusque extremae~\dots~» ; la lacune est propre à la
prop.~4, comme chez Anderson.

\paragraph{3. \emph{De emendatione}, problème~2, théorème~2 (\vue{107}).}
Le texte du manuscrit donne la constitution de l'équation
$A^4 - 2B\,A^3 = Z$ et finit sur «~et fit Z plano-plano ex applicatione
quadrati aggregati laterum minus triplo rectangulo sub lateribus ad
rectangulum [sub lateribus]~» ; ces deux lignes sont cernées d'un trait
fermé. La marge porte : «~ex applicatione quadrati \emph{cubi} a Dp. ad
differentiam [mot biffé] quadruplati quadratorum a D et B~», où «~cubi~»
est écrit au-dessus de la ligne. Le passage cerné est, mot pour mot, la
leçon qu'Anderson attribue à son exemplaire ; la note marginale est la leçon
qu'il imprime (p.~98), à l'ajout près de «~cubi~» au-dessus de
«~quadrati~», qui n'est pas biffé, et à «~quadruplati~» pour
«~quadruplam~». Le dernier mot du passage cerné est lu «~subcubis~» ; la
leçon d'Anderson est «~sub lateribus~».

\paragraph{4. \emph{De emendatione}, problème~2, théorème~3 (\vue{108}).}
Le texte du manuscrit donne la constitution de l'équation
$2B\,A^3 - A^4 = Z$ sans dire ce qu'est $Z$ ; il porte à cet endroit la
ligne «~et est A latus maius minusve~», barrée. La marge, soulignée et
reliée au texte par un trait, la remplace par «~et fit Z plano-plano ex
applicatione cubi a D pl. ad aggregatum \emph{quadruplum} quadratorum a
Bq${}+{}$Dq~», «~quadruplum~» étant ajouté au-dessus de la ligne. C'est la
phrase qu'Anderson dit absente de son exemplaire et qu'il imprime p.~99.

\subsection{Les fautes communes}

Quatre articles de la liste, sur trois pages, corrigent des fautes que le
texte du manuscrit présente aussi :
\begin{itemize}
\item p.~26, l.~2, «~addatur ipsi Ec negationis nota~» : à la \vue{32}, le
cube retranché $-E^3 + 3BE^2 - 3B^2E + B^3$ est écrit sans le signe moins
devant «~E cubo~» ; la marge ne le restitue pas ;
\item p.~44, «~pro, \emph{hisce omnibus per A diuisis} lege \emph{per
A~$-$~E}~» : la \vue{50} porte «~et omnibus per A divisis~», alors que le
diviseur écrit dessous est $A - E$ ; la marge ne corrige pas ;
\item p.~129, prop.~3, «~prioris aequationis partis fini adijciatur
$-$~Aqq~», et «~in notis Arithmeticis pro 16~C. lege 10~C.~» : à la
\vue{140}, l'énoncé de la prop.~3 du chapitre~XIV omet le terme $-A^4$ et
l'exemple écrit $16C$ ; la marge porte «~$-$~Aqq.~» et «~10~».
\end{itemize}
La faute de l'exemple est réelle :
$(x-1)(x-2)(x-3)(x-4) = x^4 - 10x^3 + 35x^2 - 50x + 24$.

La faute que cite Tannery, «~vocum~» pour «~vitium~», est hors de la
liste. L'édition de 1615 l'a aussi, dans le sommaire en italique qui suit
le titre «~De Anastrophe~» (p.~79) : «~Aut quemadmodum aduersus vocum
\emph{amphibolias}~\dots~» (le dernier mot en grec), là où le sens
demande «~adversus vitium~», comme dans les titres des chapitres~II, IV
et~V du manuscrit. Au même endroit (\vue{83}), la transcription lit
«~Et quemadmodum ad nostras \emph{vicinas}~» suivi d'un mot grec lu
\emph{hyperbolas}, avec «~vicinum~» en marge ; le mot douteux est
souligné et surligné. Cette lecture ne s'accorde ni avec «~vocum~» ni avec
«~vitium~», et le mot grec qu'elle donne n'est pas celui de l'imprimé ;
elle est sans doute fautive. Le passage est à relire sur l'image : s'il
porte «~vocum~», c'est une faute commune de plus, et la plus nette, parce
qu'elle ne peut guère venir de deux lecteurs indépendants.

\subsection{Ce que l'imprimé suit sans le dire}

La liste ne recense pas toutes les corrections de l'imprimé ; Anderson dit
d'ailleurs avoir revu tous les exemples. Sur les pages vérifiées, plusieurs
corrections de la marge du manuscrit sont dans l'imprimé sans figurer dans
la liste. À la \vue{65}, l'exemple de la prop.~5 du chapitre~19 écrit
$4\,618\,812$, nombre souligné, et la marge donne $5\,153\,632$ ; l'imprimé
(p.~61) porte «~5,153,632~». À la \vue{108}, l'exemple
«~1~QQ~$-$~5~C æquatur 216~» a son signe corrigé en «~$+$~» dans la marge, et
«~540~» y devient «~5400~» ; l'imprimé (p.~99) porte «~1~QQ.~$+$~5~C.~» et
«~5,400~». Au théorème~3, l'équation de la résolvante, absente du texte du
manuscrit, est donnée dans sa marge, et l'imprimé l'a.

\begin{table}[ht]
\centering\small
\begin{tabular}{@{}>{\raggedright}p{2.6cm}>{\raggedright}p{2.1cm}>{\raggedright}p{4.1cm}>{\raggedright\arraybackslash}p{4.6cm}@{}}
\toprule
Passage & Anderson 1615 & Texte de NAL 1644 & Marge de NAL 1644\\
\midrule
\emph{Rec.} 18, prop.~14 & p.~58 : manquait & $Z$ non défini (\vue{62}) & définition de $Z$, biffée, lue en partie\\
\emph{Rec.} 19, prop.~4 & p.~60 : manquait & phrase sans verbe, croix (\vue{64}) & «~differentia est B et fit Z\dots~» à la croix\\
\emph{Em.} pb.~2, th.~2 & p.~98 : leçon changée & leçon de l'exemplaire, cernée (\vue{107}) & leçon imprimée, «~cubi~» ajouté\\
\emph{Em.} pb.~2, th.~3 & p.~99 : manquait & ligne barrée, $Z$ non défini (\vue{108}) & «~et fit Z\dots\ quadruplum\dots~»\\
\midrule
p.~26, signe de $E^3$ & erratum & faute (\vue{32}) & non corrigée\\
p.~44, «~per A~» & erratum & faute (\vue{50}) & non corrigée\\
p.~129, $-A^4$ & erratum & omis (\vue{140}) & «~$-$~Aqq.~»\\
p.~129, $16C$ & erratum & $16$ (\vue{140}) & «~10~»\\
\bottomrule
\end{tabular}
\caption{Les articles de la liste d'Anderson qui nomment l'exemplaire, et
les errata dont la faute est aussi dans le texte du manuscrit. Lectures du
manuscrit d'après une transcription non vérifiée par un humain.}
\label{tab:collation}
\end{table}

\section{Ce qui s'accorde mal}\label{sec:desaccords}

Plusieurs articles de la liste corrigent des fautes que le texte du
manuscrit n'a pas. Les mots omis p.~15 («~quadratice, cubice, \& ulterius
Climactice. Descensus contra~\dots~») sont dans le texte de la \vue{20}, et
«~simpliciores~» (p.~16) aussi. À la p.~99, l'imprimé écrit
«~fit 1~N.~16~» et la liste corrige en «~6~» ; le manuscrit porte «~6~»
(\vue{108}). La collation faite sur les transcriptions en relève d'autres
aux p.~8, 81, 85 et~89. Ces articles relèvent sans doute de la seconde
catégorie du titre, les fautes des typographes : ils ne s'opposent à aucune
des hypothèses.

Un article est plus gênant. À la p.~112, la liste demande qu'on insère, dans
les exemples du corollaire de la formule de Cardan, «~\& 378~$-$~351 est 27,
cubus a latere 3~». Le texte de la \vue{122} a cette ligne, son «~et~»
surchargé. Si Anderson l'a ajoutée parce que son exemplaire ne l'avait pas,
le manuscrit n'est pas cet exemplaire tel qu'il l'a lu. Mais la liste ne dit
pas ici «~deerant in exemplari~» : la ligne peut avoir été sautée par le
compositeur.

L'imprimé ne suit pas non plus toutes les marges. À la \vue{108}, il garde
la ligne barrée «~et est A latus maius minusve~» après la phrase suppléée,
et la phrase «~In serie quatuor continue proportionalium B bis est composita
ex illis omnibus~», que le manuscrit barre d'un trait léger ; il ignore la
correction marginale «~$+$~Dq~» de la même page ; il ne corrige pas les
fautes de la p.~129 que la marge de la \vue{140} corrige, et doit les
reprendre dans ses errata. Les marges du manuscrit ne sont donc pas une
collation mécanique de l'imprimé, et l'imprimé n'est pas une copie
mécanique des marges.

\section{Trois hypothèses}\label{sec:hypotheses}

La lecture de Tannery tient en trois points : le manuscrit est une copie,
qui n'est ni de la main de Viète ni de celle d'Anderson ; il représente la
rédaction définitive d'Anderson ; quelques marges sont de la main
d'Anderson. La collation s'accorde avec le troisième point et nuance le
deuxième. Aux quatre endroits que la liste rapporte à l'exemplaire, le
texte du manuscrit est celui de l'exemplaire, lacunes et leçon fautive
comprises, et la rédaction définitive n'est que dans les marges. Le
manuscrit serait donc moins une copie de cette rédaction que le support
sur lequel elle a été faite, ou la copie d'un tel support. Trois
hypothèses précisent cette alternative ; la première est la plus
proche de Tannery, si l'on donne aux marges des quatre passages la main
qu'il croyait reconnaître.

\paragraph{(a) NAL 1644 est l'exemplaire d'Anderson.}
Les marges des quatre passages seraient alors la préparation de l'édition,
par Anderson lui-même ou par Aleaume, et postérieures à l'état que la liste
appelle «~l'exemplaire~». Cette hypothèse rend compte de l'emplacement des
additions, des corrections de calcul que l'imprimé suit sans les lister, de
la correction en cours d'écriture de la \vue{107} («~quadrati~» puis
«~cubi~»), et des fautes communes (p.~26, 44) que le compositeur aurait
reproduites. Elle doit expliquer les marges que l'imprimé ne suit pas
(\vue{108}, \vue{140}) : elles peuvent être postérieures à la composition.
Elle prédit que la main des quatre notes marginales est celle d'Anderson ou
d'Aleaume, et que le papier est antérieur à 1615.

\paragraph{(b) NAL 1644 est une copie de l'exemplaire, annotée plus tard.}
Un copiste aurait reproduit l'exemplaire d'Anderson, lacunes comprises, puis
les additions et corrections portées dans ses marges, ou l'imprimé. Les
notes «~transcripta ex exemplari~» et «~deerant haec in exemplari~», si
elles se lisent ainsi, vont dans ce sens, comme l'impression que les marges
des \vues{105}{119} sont de la main et de l'encre du texte. On attendrait
alors une écriture de copiste dans le texte et les marges, éventuellement
un papier postérieur à 1615. Les marges que Tannery croyait de la main
d'Anderson, si elles le sont, s'expliquent mal dans ce cas.

\paragraph{(c) NAL 1644 dérive de la même source et a été corrigé d'après l'imprimé.}
Le texte viendrait d'un autre témoin des papiers de Viète, défectueux aux
mêmes endroits, et un lecteur l'aurait complété sur l'édition de 1615 et sa
liste. Les écarts de formulation («~utriusque inter extremas~» contre
«~utriusuis extremae~», «~quadruplati~» contre «~quadruplam~») et les marges
que l'imprimé ne suit pas, ou qu'il contredit (lignes barrées qu'il garde),
s'accordent mal avec une simple collation sur le livre ; ils ne
l'excluent pas.

\paragraph{Ce qui trancherait.}
La comparaison des écritures du texte et des marges des \vues{61}{66} et
\vues{105}{108}, et des feuillets~A et~B (\vues{2}{3}), avec des
autographes d'Anderson et d'Aleaume, et avec le lat.~7274 que Bigourdan
rapproche de la main principale ; pour Anderson, l'échantillon est celui
que désignait Tannery, les notes de sa main sur l'exemplaire du Diophante
de Bachet conservé à la Sorbonne ; la relecture de la \vue{83}, où
«~vocum~» ferait une faute commune de plus ; l'examen du papier
(filigranes) et des encres sur l'original ; la collation complète du
manuscrit avec l'édition de 1615, qui dirait si d'autres corrections
silencieuses d'Anderson sont dans ses marges ; la lecture assurée des notes
des vues~20, 100 et~114. Si le texte était de la main de Viète, comme
Bigourdan l'envisage pour les feuillets~1 à~66, l'hypothèse~(c) perdrait
beaucoup et l'exemplaire d'Anderson pourrait être un autographe.

\section{Les énoncés en notation moderne}\label{sec:moderne}

Les quatre passages portent sur des équations dont Viète donne la
«~constitution~», c'est-à-dire les racines en fonction de grandeurs en
proportion continue. En notation moderne, et pour des grandeurs positives :

\begin{proposition}[\emph{Rec.}~18, prop.~14]\label{prop:rec18}
Soient $a, b, c$ en proportion continue ($b^2 = ac$), $s = a + b$,
$t = b + c$, $B = s^3 + t^3 + a\,t^2$ et $Z = (B - s^3)\,s^2$. Alors
$s$ et $t$ sont racines de $B x^2 - x^5 = Z$.
\end{proposition}

\noindent Pour $s$ c'est la définition de $Z$, qui est le passage suppléé
par Anderson ; pour $t$, on vérifie $(B - s^3)s^2 = (B - t^3)t^2$ en posant
$a = 1$, $b = r$, $c = r^2$ : les deux membres valent $r^2 s^4(1 + r + r^2)$.
Sur $1, 2, 4$, on trouve $B = 279$, $Z = 2268$, et les racines $3$ et $6$ de
l'exemple du manuscrit et de l'imprimé.

\begin{proposition}[\emph{Rec.}~19, prop.~4]\label{prop:rec19}
Soient $a, b, c, d$ en proportion continue, $a < d$,
$B = (b + d) - (a + c)$ et $Z = a(b - c + d)^3 = d(a - b + c)^3$. Alors
$x = a - b + c$ est racine de $x^4 + Bx^3 = Z$ et $y = b - c + d$ est racine
de $y^4 - By^3 = Z$.
\end{proposition}

\noindent En effet $x + B = d$ et $y - B = a$, d'où $x^3(x + B) = d\,x^3 = Z$
et $y^3(y - B) = a\,y^3 = Z$. La phrase que l'exemplaire n'avait pas est la
définition de $B$ et de $Z$. Sur $1, 2, 4, 8$ : $B = 5$, $Z = 216$, $x = 3$,
$y = 6$.

\begin{proposition}[\emph{Em.}, pb.~2, th.~2 et~3]\label{prop:em}
La quartique $x^4 - 2Bx^3 = Z$ se ramène au second degré par la résolvante
$D^3 + 4ZD = 4ZB^2$, et la quartique $2Bx^3 - x^4 = Z$ par
$D^3 - 4ZD = 4ZB^2$. Ainsi $Z = D^3 / \bigl(4(B^2 - D)\bigr)$ dans le
premier cas, et $Z = D^3 / \bigl(4(B^2 + D)\bigr)$ dans le second.
\end{proposition}

\noindent C'est la méthode qu'on attribue aujourd'hui à Ferrari : on forme
le carré de $x^2 - Bx - \tfrac12 D$ et l'on choisit $D$ pour que l'autre
membre soit aussi un carré. Les deux phrases d'Anderson, «~ex applicatione
cubi a D plano ad differentiam [aggregatum] quadruplam quadratorum a D \& B~»,
ne sont que la résolvante résolue en~$Z$. La leçon de l'exemplaire est
fausse : sur l'exemple $1, 2, 4, 8$ du manuscrit et de l'imprimé, les deux
côtés sont $3$ et $6$ ($2B = (6^4 - 3^4)/(6^3 + 3^3) = 5$), et la leçon
cernée donne $\bigl((3+6)^2 - 3\cdot 18\bigr)/18 = 3/2$, quand $Z = 216$ ;
elle n'a pas non plus la dimension d'un \emph{plano-planum}. La leçon
d'Anderson donne, avec $D = 6$ (racine de $D^3 + 864D = 5400$, l'équation
«~$1C + 864N$ æquatur 5400~» de l'exemple), $6^3 / \bigl(4(6{,}25 - 6)\bigr)
= 216$.

\section*{Questions ouvertes}

Plusieurs travaux n'ont pas été lus et peuvent déjà contenir tout ou
partie de ce qui précède. L'article de Busard sur les papiers de Viète à la
Bibliothèque nationale \cite{busard1964} porte surtout sur le NAL 1643,
mais touche au NAL 1644 et en dit les ff.~67--79 autographes, d'après une
source secondaire ; il faut savoir ce qu'il dit des traités algébriques et
de leurs marges. La thèse de Jean Grisard \cite{grisard1970}, dont le seul
exemplaire signalé est à la réserve de la Bibliothèque de la Sorbonne, est
le premier endroit où une collation avec l'édition de 1615 pourrait se
trouver. Il faut lire en entier la notice du \emph{De recognitione et
emendatione} dans le catalogue de Van~Egmond (n\textsuperscript{o}~15)
\cite{vanegmond1985}, et les pages de Hofmann qui mentionnent le
lat.~1644 à propos de ces deux traités \cite{hofmann1962}. La traduction
de Witmer \cite{witmer1983} ne paraît pas traiter du manuscrit, ce qui n'a
pas été vérifié page à page. La comparaison des mains suppose des
échantillons de l'écriture d'Anderson et d'Aleaume ; pour Anderson, le
Diophante annoté de la Sorbonne que citait Tannery est le premier à
consulter. Toutes les lectures du manuscrit citées ici restent à
contrôler sur les images, en particulier la note biffée de la vue~62,
«~utriusque~» à la vue~64, le cerne de la vue~107, les traits qui barrent
les lignes de la vue~108 et le sommaire du chapitre~III de la vue~83.

\section*{Sources et vérification}

Les lectures de NAL 1644 viennent d'une première transcription du volume,
faite vue par vue sur les images de Gallica par un modèle de langue et
publiée avec son apparat critique sur \url{https://germain.commutator.io}
(par exemple \url{https://germain.commutator.io/archive/#nal-1644/4/fr/62}).
Elle n'a été vérifiée sur les images par aucun humain, et aucune des lectures
du manuscrit données ici n'a été recontrôlée sur l'image pour cet article.
Les citations de l'édition de 1615 (préface, liste, p.~58, 60, 61, 98
et~99) ont été lues sur les images de la numérisation d'archive.org
\cite{anderson1615}, ainsi que la p.~79 ; les citations de Bigourdan et de
Delisle, sur l'OCR des numérisations d'archive.org ; la question de Braid et
la réponse de Tannery, sur les images de la réimpression des
\emph{Mémoires scientifiques} \cite{tannery1902}. La recension de 1896
\cite{tannery1896}, le catalogue de Van~Egmond \cite{vanegmond1985} et
l'article de Busard \cite{busard1964} ne sont connus ici que par des
extraits ou de seconde main, et Grisard, Hofmann et Witmer n'ont pas été
lus. Les calculs des propositions~\ref{prop:rec18}
à~\ref{prop:em} ont été refaits à la main.

La transcription, les recherches bibliographiques et la rédaction de
l'article ont été réalisées avec l'aide d'un modèle de langue (Claude,
d'Anthropic), sous la direction de l'auteur.

{\small\begin{thebibliography}{99}\setlength{\itemsep}{1pt}

\bibitem{anderson1615}
F.~Viète, \emph{De aequationum recognitione et emendatione tractatus duo},
éd. A.~Anderson, Paris, J.~Laquehay, 1615.
\url{https://archive.org/details/bub_gb_Q4dbiZ0apcUC}. La préface
«~Ad Matheseos studiosos~» et la liste des suppléments occupent les
quatre pages non paginées qui précèdent le texte (images n12 à n15 de la
numérisation).

\bibitem{bigourdan1916}
G.~Bigourdan, «~Sur un ouvrage de F.~Viète, supposé perdu :
\emph{Harmonicon coeleste}~», \emph{Comptes rendus hebdomadaires des
séances de l'Académie des sciences} 162 (1916), p.~237--240.
\url{https://archive.org/details/comptes-rendus-hebdomadaires-academie-des-sciences_1916-02-14_162_7}.

\bibitem{braid1902}
H.~Braid, question 2406, \emph{L'Intermédiaire des mathématiciens} 9
(1902), p.~206 ; réimprimée avec la réponse de Tannery dans
\cite[p.~410]{tannery1902}.

\bibitem{busard1964}
H.~L.~L.~Busard, «~Über einige Papiere aus Viètes Nachlaß in der Pariser
Bibliothèque Nationale~», \emph{Centaurus} 10 (1964), p.~65--126.
\href{https://doi.org/10.1111/j.1600-0498.1964.tb00393.x}{doi:10.1111/j.1600-0498.1964.tb00393.x}.

\bibitem{delisle1888}
L.~Delisle, \emph{Catalogue des manuscrits des fonds Libri et Barrois},
Paris, H.~Champion, 1888.
\url{https://archive.org/details/cataloguelibri00bibl}.

\bibitem{delisle1891}
L.~Delisle, \emph{Manuscrits latins et français ajoutés aux fonds des
nouvelles acquisitions pendant les années 1875--1891}, t.~II, Paris, 1891.

\bibitem{gallica}
Bibliothèque nationale de France, Nouvelles acquisitions latines 1644,
\url{https://gallica.bnf.fr/ark:/12148/btv1b10033620q}.

\bibitem{grisard1970}
J.~Grisard, \emph{François Viète, mathématicien de la fin du seizième
siècle : essai bio-bibliographique}, thèse de 3\textsuperscript{e} cycle,
Lettres, Paris, 1970 (date du SUDOC, PPN 174496044 ; la thèse est datée de
1968 par Van~Egmond et par un recenseur de Witmer). Bibliothèque de la
Sorbonne, Réserve I~4=~1248 (1--2).

\bibitem{henry1880}
C.~Henry, \emph{Recherches sur les manuscrits de Pierre de Fermat}, 1880.

\bibitem{hofmann1962}
J.~E.~Hofmann, «~Über Viètes Beiträge zur Geometrie der Einschiebungen~»,
\emph{Mathematisch-Physikalische Semesterberichte} 8 (1962), p.~191--214.

\bibitem{opera1646}
F.~Viète, \emph{Opera mathematica}, éd. F.~van Schooten, Leyde,
B. et A.~Elzevier, 1646.
\url{https://archive.org/details/bub_gb_25iHrtJpy9oC}.

\bibitem{pantin2003}
I.~Pantin, «~Les \emph{Principes de cosmographie} (1637). Quelques
interrogations~», \emph{Littérales}, hors-série 3 (2003), p.~117--133.
\href{https://doi.org/10.3406/litra.2003.1906}{doi:10.3406/litra.2003.1906}.

\bibitem{tannery1896}
P.~Tannery, compte rendu de F.~Ritter, \emph{François Viète, inventeur de
l'algèbre moderne}, \emph{Bulletin des sciences mathématiques},
2\textsuperscript{e} sér., 20, 1\textsuperscript{re} partie (1896),
p.~204--211. HathiTrust hvd.32044102908712.

\bibitem{tannery1902}
P.~Tannery, réponse à la question 2406, \emph{L'Intermédiaire des
mathématiciens} 9 (1902), p.~329 ; réimpr. dans P.~Tannery,
\emph{Mémoires scientifiques}, t.~X, 1930,
p.~410--411, n\textsuperscript{o}~21.
\url{https://archive.org/details/in.ernet.dli.2015.377303}.

\bibitem{vanegmond1985}
W.~Van~Egmond, «~A catalog of François Viète's printed and manuscript
works~», dans M.~Folkerts et U.~Lindgren (éd.), \emph{Mathemata. Festschrift
für Helmuth Gericke}, Stuttgart, F.~Steiner, coll. «~Boethius~» 12, 1985,
p.~359--396.

\bibitem{witmer1983}
F.~Viète, \emph{The Analytic Art: Nine Studies in Algebra, Geometry and
Trigonometry}\dots, trad. T.~R.~Witmer, Kent (Ohio), Kent State University
Press, 1983, 450~p.

\end{thebibliography}}
\end{document}
